\documentclass[spanish]{article}
\def\spanishoptions{mexico}
\usepackage[spanish]{babel}      % Para procesar adecuadamente el español
\usepackage[utf8]{inputenc}      % Codificación de entrada UTF-8
\usepackage[T1]{fontenc}         % Codificación de salida T1
\usepackage[official]{eurosym}   % Símbolo de euro
\usepackage{siunitx}             % Paquete para unidades métricas
\usepackage{textcomp}            % Tipografías compañeras
\usepackage{lmodern}             % Tipografía Latin Modern
\usepackage{moreverb}            % Para escribir código
\usepackage{tfrupee}
\usepackage{fontawesome5}
\usepackage{parskip}
\usepackage{hyperref}
\usepackage{tasks}
\usepackage{exsheets}
\usepackage{enumitem}
\usepackage{lastpage}
\usepackage{booktabs}
\usepackage{graphicx}
\usepackage{apalike}
\usepackage{relsize}
\usepackage{listings}
\usepackage{amsmath}
\usepackage{amssymb}
\usepackage{fontawesome5}

% Formato global:
\setlength{\parindent}{0cm}      % Párrafos sin sangría inicial
\hypersetup{pdfborder={0 0 0}}   % Ligas sin borde

\begin{document}
\title{Proyecto final de Optimización Convexa \\ Optimización de hiperparámetros en modelos predictivos}
\author{Acoyani Garrido Sandoval}
\date{\today}
\maketitle




\section{Introducción}

El objetivo del presente proyecto final es aplicar los conceptos aprendidos en el curso de Optimización Convexa para resolver un problema de optimización de hiperparámetros en un modelo predictivo.

En general, los pasos del proyecto son:

\begin{enumerate}
   \item Seleccionar un problema de aprendizaje supervisado o no supervisado. Puede ser uno que esté siendo desarrollado, o que tenga relevancia personal o académica.
   \item Elegir modelos base que tengan hiperparámetros susceptibles de ser optimizados.
   \item Identificar los hiperparámetros críticos del modelo y definir un espacio de búsqueda adecuado para cada uno de ellos.
   \item Plantear el problema de buscar los valores de los hiperparámetros como un problema de optimización: regularización L1, L2 o \textit{elasticnet}, conjuntos convexos y restricciones, minimización de funciones convexas, dualidad y condiciones de optimalidad, o métodos de optimización tales como gradiente descendente.
   \item Implementar la solución del problema en una Jupyter notebook. Está permitido el uso de librerías, siempre y cuando se justifique el método de optimización.
\end{enumerate}




\clearpage
\section{Selección de conjunto de datos}

El conjunto de datos usado para estos fines es el \textbf{Telco Customer Churn} \cite{kaggletelco2026}, un conjunto de datos desarrollado por IBM originalmente como una muestra de datos incluída con el software \textbf{IBM Cognos Analytics Server} \cite{IBMcognos2026}, un software de análisis y visualización de datos basado en Java Enterprise Edition. Este conjunto de datos reviste una doble importancia personal para mí, ya que por una parte se trata de un conjunto de datos relacionado con comportamiento de clientes, giro en el cual actualmente manejo un negocio denominado \textbf{BuzónQR\textregistered{}}; y por otro lado, por el hecho de ser un conjunto de datos inicialmente concebido para probar las capacidades del software IBM Cognos Analytics Server, una plataforma de inteligencia de negocios inicialmente desarrollada por la empresa canadiense Cognos que fue comprada por IBM hacia el año 2008, orientada principalmente a visualización, procesamiento, análisis y exploración de datos; mi primer trabajo, el cual desempeñé en IBM desde 2012 hasta 2018, fue como administrador de una serie de instancias de dicha plataforma.



\subsection{Características del conjunto de datos}

El conjunto de datos \textit{Telco Customer Churn} tiene como finalidad ofrecer un medio para mejorar la retención de clientes a través de información acerca de ellos y un dato que muestra si han dejado la empresa de telefonía en el último mes.

En este conjunto de datos, cada fila representa uno de los usuarios de la empresa de telefonía, y cada columna representa información acerca del cliente. Las características de este modelo son:

\begin{enumerate}
   \item \textbf{Variable dependiente:} \textit{Churn}. Es una variable binaria cuyos valores son ``Yes'' o ``No'', e indica si el cliente ha dejado la compañía en el último mes. La distribución de clases de esta variable es de 73.5\% retenidos ($\mathrm{Pr}(Churn = No)$) y 26.5\% no retenidos ($\mathrm{Pr}(Churn = Yes)$); está moderadamente desbalanceado, lo cual es importante a la hora de elegir hiperparámetros de ponderación de clase.
   \item \textbf{Características:} Encuadran en 4 grupos naturales: \textbf{demográficas} (\textit{gender, SeniorCitizen, Partner, Dependents}), \textbf{información de cuenta} (\textit{tenure} (meses como cliente), \textit{Contract} (mensual, anual, bienal), \textit{PaperlessBilling, PaymentMethod}), y \textbf{servicios suscritos} (\textit{PhoneService, MultipleLines, InternetService, OnlineSecurity, OnlineBackup, DeviceProtection, TechSupport, StreamingTV, StreamingMovies}), y \textbf{cargos} (\textit{MonthlyCharges, TotalCharges}).
\end{enumerate}

Si realizamos un proceso de \textit{one-hot encoding} de las características (descrito más adelante), obtenemos un total de 30 dimensiones de datos; lo cual implica que \textbf{será necesario elegir características a través de la regularización}, a la vez que la cantidad de características es lo suficientemente pequeña para que un solver convexo lo pueda resolver.

En general, podemos observar que se trata de un problema de clasificación supervisada binaria, lo cual hace que el modelo natural a elegir sea \textbf{regresión logística}. Además de prestarse de forma natural a variables dependientes binarias, este modelo tiene 3 características que van en línea con el contenido del curso:

\begin{enumerate}
   \item La regresión logística tiene una función de log-verosimilitud negativa convexa en parámetros, lo que hace que la superficie de pérdida tenga un mínimo global único.
   \item Añadir regularización L1, L2 o elasticnet preservará la convexidad, lo que nos permitirá aplicarlas sin complejidades.
   \item El problema de optimización tiene un planteamiento dual y condiciones KKT bien conocidas.
\end{enumerate}



\subsection{One-hot encoding de características}

En la subsección anterior, mencioné que un proceso de \textit{one-hot encoding} de las características arrojó 30 dimensiones de datos. El proceso para hacer eso y las implicaciones se explican a continuación.

\paragraph{¿Qué es?} El \textit{one-hot encoding} (codificación en caliente) es una técnica de preprocesamiento que convierte variables categóricas en variables numéricas binarias. Para cada variable categórica con $k$ categorías distintas, se crean $k$ columnas binarias (indicadoras), una por categoría, donde el valor 1 indica que la observación pertenece a esa categoría y 0 en caso contrario. En la práctica, para evitar multicolinealidad perfecta (\textit{dummy variable trap}), se elimina una de las $k$ columnas por variable, resultando en $k-1$ columnas nuevas; la categoría eliminada queda implícitamente representada cuando todas las demás son 0.

\paragraph{Proceso} El conjunto de datos \textit{Telco Customer Churn} contiene 19 características (excluyendo \textit{customerID} y la variable dependiente \textit{Churn}). De ellas, 3 son ya numéricas continuas (\textit{tenure, MonthlyCharges, TotalCharges}) y 1 es binaria numérica (\textit{SeniorCitizen}), por lo que se mantienen sin transformación. Las 15 características categóricas restantes se procesan con \texttt{pd.get\_dummies(..., drop\_first=True)} de la librería \textit{Pandas}, eliminando una categoría por variable para evitar la trampa de la variable ficticia. La tabla siguiente resume la expansión de dimensiones:

\begin{center}
\begin{tabular}{llr}
\toprule
\textbf{Característica} & \textbf{Categorías} & \textbf{Columnas} \\
  &   & \textbf{generadas} \\
\midrule
\textit{gender}           & Female, Male                            & 1 \\
\textit{Partner}          & No, Yes                                 & 1 \\
\textit{Dependents}       & No, Yes                                 & 1 \\
\textit{PhoneService}     & No, Yes                                 & 1 \\
\textit{PaperlessBilling} & No, Yes                                 & 1 \\
\textit{MultipleLines}    & No, No phone service, Yes               & 2 \\
\textit{InternetService}  & DSL, Fiber optic, No                    & 2 \\
\textit{OnlineSecurity}   & No, No internet service, Yes            & 2 \\
\textit{OnlineBackup}     & No, No internet service, Yes            & 2 \\
\textit{DeviceProtection} & No, No internet service, Yes            & 2 \\
\textit{TechSupport}      & No, No internet service, Yes            & 2 \\
\textit{StreamingTV}      & No, No internet service, Yes            & 2 \\
\textit{StreamingMovies}  & No, No internet service, Yes            & 2 \\
\textit{Contract}         & Month-to-month, One year, Two year      & 2 \\
\textit{PaymentMethod}    & Bank transfer, Credit card,  & 3 \\
                          & Electronic check, Mailed check & 3 \\
\midrule
4 numéricas               & (\textit{SeniorCitizen, tenure,}) & 4 \\
                          & (\textit{MonthlyCharges, TotalCharges}) & 4 \\
\midrule
\textbf{Total}            &                                         & \textbf{30} \\
\bottomrule
\end{tabular}
\end{center}

\paragraph{Implicaciones}

\begin{enumerate}
   \item \textbf{Dimensionalidad manejable.} Las 30 dimensiones resultantes son suficientemente pocas para que un solver convexo (como \texttt{liblinear} o \texttt{lbfgs} dentro de \texttt{sklearn.linear\_model.LogisticRegression}) las maneje eficientemente.
   \item \textbf{Necesidad de regularización.} Con 30 características y 7{,}043 observaciones, la razón de observaciones por característica es de aproximadamente 235:1, lo cual en principio es favorable. Sin embargo, varios subgrupos de variables indicadoras están correlacionados entre sí (p.~ej., los servicios que dependen de \textit{InternetService}), por lo que la regularización L1 o L2 sigue siendo conveniente para controlar la varianza del modelo y producir coeficientes interpretables.
   \item \textbf{Interpretabilidad de coeficientes.} Cada coeficiente $\beta_j$ en la regresión logística representa el cambio en el log-odds de \textit{Churn} al pasar de la categoría de referencia a la categoría representada, manteniendo las demás características constantes; esto facilita la interpretación de resultados.
   \item \textbf{Escala.} Las columnas numéricas continuas (\textit{tenure, MonthlyCharges, TotalCharges}) tienen escalas distintas a las variables indicadoras $\{0,1\}$, por lo que es necesario estandarizarlas (media cero, varianza unitaria) antes de ajustar el modelo con regularización, para que la penalización afecte a todos los coeficientes de forma equitativa.
\end{enumerate}




\section{Desarrollo del problema}

Habiendo expuesto la composición de nuestro conjunto de datos, la metodología que seguiremos será:

\begin{enumerate}
   \item \textbf{Planteamiento del problema:} Predecir $\mathrm{Pr}(Churn = 1 | x)$ a partir de las características de los clientes.
   \item \textbf{Modelos:} Regresión logística como modelo primario, SVM lineal como modelo secundario. Compararé la pérdida logística contra la \textit{hinge loss}. Privilegiaré el uso de modelos lineales con el fin de mantenernos en un marco de referencia convexo.
   \item Plantear los hiperparámetros a optimizar y su espacio de búsqueda.
   \item Plantear un problema de optimización convexa.
   \item Implementación.
\end{enumerate}



\clearpage
\clearpage
\section{Modelo predictivo de regresión logística}



\bibliographystyle{apalike}
\bibliography{../bibliografias}

\end{document}
